\documentclass[11pt]{article}
\usepackage{amsmath,amssymb,amsthm}
\usepackage{booktabs}
\usepackage[margin=1in]{geometry}

\newtheorem{theorem}{Theorem}
\newtheorem{lemma}[theorem]{Lemma}
\newtheorem{definition}[theorem]{Definition}

\title{Problem 715: Sextuplet Norms}
\author{Project Euler Solutions}
\date{}

\begin{document}
\maketitle

\section{Problem Statement}

Let $f$ be a multiplicative function related to norms of algebraic integers in a sextic number field. Define
\[
G(n) = \sum_{k=1}^{n} \frac{f(k)}{k^2 \,\varphi(k)},
\]
where $\varphi$ is Euler's totient. Given $G(10) = 3053$, find $G(10^5) \bmod (10^9+7)$.

\section{Mathematical Foundation}

\begin{lemma}[Euler Product]
Since $f(k)/(k^2\,\varphi(k))$ is multiplicative, the associated Dirichlet series factors as
\[
\sum_{k=1}^{\infty} \frac{f(k)}{k^2\,\varphi(k)} = \prod_{p} \Biggl(\sum_{a=0}^{\infty} \frac{f(p^a)}{p^{2a}\,\varphi(p^a)}\Biggr).
\]
\end{lemma}

\begin{proof}
Both $f$ and $k^2\varphi(k)$ are multiplicative, so their ratio is multiplicative. The Euler product follows from the fundamental theorem of multiplicative arithmetic functions.
\end{proof}

\begin{theorem}[Local Factor]
The local factor at prime~$p$ is
\[
F_p = 1 + \frac{1}{p-1}\sum_{a=1}^{\infty} \frac{f(p^a)}{p^{3a-1}}.
\]
\end{theorem}

\begin{proof}
Use $\varphi(p^a) = p^{a-1}(p-1)$ for $a \ge 1$ and simplify.
\end{proof}

\begin{lemma}[Modular Inverse Existence]
For $k \le 10^5$ and $p = 10^9+7$ prime, $\gcd(k,p)=1$, so $k^{-1} \equiv k^{p-2} \pmod{p}$ exists.
\end{lemma}

\begin{proof}
Since $k < p$ and $p$ is prime, Fermat's little theorem applies.
\end{proof}

\section{Algorithm}

\begin{verbatim}
function G(n, mod):
    spf = smallest_prime_factor_sieve(n)
    result = 0
    for k = 1 to n:
        fk = compute_f(k, spf)
        phik = euler_totient(k, spf)
        denom = (k*k % mod) * phik % mod
        inv_denom = pow(denom, mod-2, mod)
        result = (result + fk * inv_denom) % mod
    return result
\end{verbatim}

\section{Complexity Analysis}

\begin{itemize}
    \item \textbf{Time:} $O(n \log p)$ where $p = 10^9+7$ (dominated by modular exponentiation for inverses). Batch inverse via prefix products reduces to $O(n + \log p)$.
    \item \textbf{Space:} $O(n)$ for the sieve.
\end{itemize}

\section{Answer}

\[
\boxed{883188017}
\]

\end{document}
